% titel: Blatt 2, A3: Stoffmenge und Teilchenzahl in 42 g
% skript: Mortimer Kap. 4.2 (Mol, Stoffmenge, Avogadro)

% Teil I: Werkzeuge, kein Lösungsweg.

Zwischen Masse, Stoffmenge und Teilchenzahl einer Verbindung vermitteln die
molare Masse $M$ und die Avogadro-Konstante $N_\mathrm{A}$. Die Aufgabe verlangt
beide Umrechnungen für eine gegebene Masse.

\begin{werkzeugkasten}[Masse --- Stoffmenge --- Teilchenzahl]
\begin{equation}
  n = \frac{m}{M}, \qquad N = n \cdot N_\mathrm{A}
  = \frac{m}{M}\cdot N_\mathrm{A},
  \qquad N_\mathrm{A} = \qty{6.022}{e23 mol^{-1}}.
\end{equation}
$n$ ist die Stoffmenge in $\mathrm{mol}$, $m$ die Masse in $\mathrm{g}$ und $M$
die (molare) Masse in $\mathrm{g/mol}$; $N$ ist die Anzahl der Teilchen (hier
Moleküle). Zwei Schritte, immer in dieser Reihenfolge: erst $n$, dann $N$.

Molmassen der Verbindungen braucht man zunaechst selbst (siehe A2):
zum Beispiel $\mathrm{HCl}$, $\mathrm{NaCl}$, Gips $\mathrm{CaSO_4}$ und
Haushaltszucker (Saccharose $\mathrm{C_{12}H_{22}O_{11}}$).
\textit{Skript:} Mortimer Kap. 4.2 (Mol, Stoffmenge, $N_\mathrm{A}$),
Übung 2/3.
\end{werkzeugkasten}

\begin{achtung}[Darauf achten]
\begin{itemize}
  \item Erst die Molmasse $M$ (Summe der Atommassen) bestimmen, dann
        einsetzen --- das ist der haeufigste Fehler.
  \item Einheiten: $m$ in $\mathrm{g}$, $M$ in $\mathrm{g/mol}$, so kürzt
        sich $\mathrm{mol}$ sauber heraus; $N$ ist eine reine Zahl (Teilchen).
  \item Potenzen bei $N_\mathrm{A}$ ($10^{23}$) und beim Ergebnis sauber
        mitschreiben; grosse Ergebnisse in der Form $x\cdot10^{23}$ angeben.
  \item Teilchenzahl abhaengig von der Verbindung; bei gleicher Masse haben
        Verbindungen mit größerer Molmasse weniger Moleküle.
\end{itemize}
\end{achtung}
